\section{Statement, representation, and logical scope}
\label{sec:scope}

\subsection{The ordinary geometric criterion}

Let \(T\subseteq\R^2\) be a closed topological disk.  Choose six consecutive boundary points
\[
 A,B,C,D,E,F
\]
in counterclockwise order; adjacent chosen points are allowed to coincide, so an intervening boundary part may be empty.  Write
\[
 a=AB,\quad b=BC,\quad c=CD,\quad d=DE,\quad e=EF,\quad f=FA
\]
for the corresponding oriented boundary paths.  Conway's condition is:
\begin{enumerate}[label=(\roman*)]
\item a translation takes \(A\) to \(E\), \(B\) to \(D\), and the path \(a\) to the path \(ED=\rev(d)\);
\item each of \(b,c,e,f\) is invariant, as an unoriented path, under the half-turn interchanging its endpoints;
\item at least three of \(A,\ldots,F\) are distinct.
\end{enumerate}
The conclusion is that congruent copies of \(T\), obtained using translations and half-turns, tile the plane periodically.  This is the labelled formulation emphasized by Schattschneider~\cite{Schattschneider1980}; it removes the ambiguity in saying merely that two boundary parts are ``parallel.''

The strip proof is geometrically persuasive: translate copies along \(a,d\) to form a strip, rotate a second strip around a midpoint of \(e\) or \(f\), and repeat.  A modern exposition explicitly observes that this description is still vague at the vertices and cites simple connectivity and covering spaces as a repair~\cite[Remark~3.4.4]{Weber2025}.  The purpose of this report is to make that repair finite and arithmetical.

\subsection{Why an arithmetic presentation is needed}

First-order iPRA has natural numbers as objects.  It does not directly quantify over arbitrary subsets of \(\R^2\), arbitrary continuous maps, or an uncountable family of tile copies.  A theorem in iPRA must therefore specify:
\begin{itemize}
\item how a tile and its boundary are number-coded;
\item which finite certificate says that the code is a disk;
\item which primitive-recursive functions produce the proposed tiling certificate; and
\item which bounded, decidable predicate verifies that certificate.
\end{itemize}

The exact theorem proved here uses a \emph{rational piecewise-linear disk}.  After multiplying all coordinates by a common denominator, one may equivalently use integer coordinates.  A disk certificate consists of a finite list of positively oriented triangles with exact edge-incidence data.  The union of the triangles is the tile, their interiors are pairwise disjoint, every interior edge occurs twice with opposite orientations, and the remaining oriented edges form one simple boundary cycle.  The six boundary paths are finite vertex lists along that cycle.  A geometrically empty path is represented by a singleton list.

This covers polygons, polyominoes, polyiamonds, polyhexes, and every exact finite straight-line model normally used in applications of the criterion.  It is also the correct finite kernel for extensions to effectively named curved boundaries: an extension must supply enough modulus and approximation data to reduce the relevant claim to compatible finite PL certificates.  We do not silently identify that stronger represented-analysis theorem with the finite theorem proved below.

\begin{definition}[Normalized PL Conway datum]
A \emph{normalized PL Conway datum} is a positively oriented rational PL disk \(T\), a certified triangulation of \(T\), and six nonempty vertex lists \(a,b,c,d,e,f\) whose concatenation is the oriented boundary cycle, satisfying
\begin{align}
 \trans{t}(a)&=\rev(d), & t&=E-A, \label{eq:ad}\\
 \half{B+C}(b)&=\rev(b), \label{eq:bhalf}\\
 \half{C+D}(c)&=\rev(c), \label{eq:chalf}\\
 \half{E+F}(e)&=\rev(e), \label{eq:ehalf}\\
 \half{F+A}(f)&=\rev(f). \label{eq:fhalf}
\end{align}
Here \(\trans{w}(x)=x+w\), while \(\half{s}(x)=s-x\) is the half-turn whose \emph{doubled center} is \(s\).  Equality of paths means equality after deleting consecutive duplicate vertices and redundant collinear subdivision vertices.
\end{definition}

Using doubled centers avoids division by two.  Every displayed equation is an equality of finite lists of rational points, hence is decidable by cross multiplication.

\subsection{The explicit theorem}

Set
\begin{equation}
  t=E-A,
  \qquad
  u=E+F-B-C,
  \qquad
  R=\half{E+F}.
  \label{eq:periods}
\end{equation}
For \((m,n)\in\Z^2\), define two tile orbits
\begin{equation}
 T^0_{m,n}=T+mt+nu,
 \qquad
 T^1_{m,n}=R(T)+mt+nu.
 \label{eq:tilefamily}
\end{equation}

\begin{theorem}[Finite PL Conway criterion]
\label{thm:mainPL}
For every normalized PL Conway datum:
\begin{enumerate}[label=(\alph*)]
\item \(\det(t,u)=2\Area(T)>0\), so \(t,u\) generate a rank-two lattice;
\item the family in \cref{eq:tilefamily} covers \(\R^2\);
\item distinct members have disjoint interiors; and
\item every boundary pairing is a translation or a half-turn generated by the four specified centrosymmetries.
\end{enumerate}
Consequently \(T\) admits a periodic monohedral tiling using only translations and half-turns.  A lattice fundamental cell contains exactly two tile copies by area.
\end{theorem}

The proof occupies \cref{sec:affine,sec:orbit,sec:area,sec:torus,sec:mainproof}.  It is constructive in a stronger sense than mere intuitionistic acceptability: from the finite datum it computes the periods, the two orbit representatives, a torus subdivision, and every bounded patch.

\subsection{The arithmetic endpoint}

Let \(c\) be a natural-number code for all finite input data.  The intended primitive-recursive predicates are
\[
 \code{DiskOK}(c),\quad \code{ConwayOK}(c),\quad
 \code{TorusOK}(c,z),\quad \code{PatchOK}(c,r,z),
\]
and the constructors are \(\code{BuildTorus}(c)\) and \(\code{BuildPatch}(c,r)\).  The central iPRA sentence has the schematic form
\begin{equation}
 \forall c\;\Bigl(
  \code{DiskOK}(c)=1\wedge\code{ConwayOK}(c)=1
  \;\to\;
  \code{TorusOK}\bigl(c,\code{BuildTorus}(c)\bigr)=1
 \Bigr).
 \label{eq:ipraschema}
\end{equation}
A companion universal formula quantifies also over a radius \(r\) and verifies the lifted finite patch.  Because every displayed function is primitive recursive, the matrix of \cref{eq:ipraschema} is quantifier-free after the usual definitional expansion.  Thus the proof uses only defining equations, equality, intuitionistic propositional reasoning, and quantifier-free induction.

\begin{remark}[What is, and is not, claimed]
The human-readable proof of \cref{thm:mainPL} is complete.  The report also gives a precise implementation plan for \cref{eq:ipraschema}.  The supplied Lean file does not yet construct the final empty-context \(\Prf\,\PRA\) derivation.  Nor do we claim a first-order quantification over all unrepresented topological disks.  These boundaries are essential to a trustworthy formalization statement.
\end{remark}
